\clearpage
\section{Full BPerm robustness by whole-graph rewrites}
\label{sec:copying-full-bperm}
\termdef{full-bperm-robustness}
Both directions of BPerm robustness are now proved for every original
recursive specimen of the concrete copying UAM. The conclusion includes an
actual rewrite of the opposite whole graph and an original residual specimen
at the actual resumed token. No matching-context, local-diagram, residual,
or robustness hypothesis remains in the final theorems.

This closes the BPerm part of the paper's Table~3 and
\href{https://arxiv.org/html/1907.01257v5\#A8.SS3.SSS6}{Appendix H.3.6}.
The proof follows its primitive, opening and copying overlap distinction.
It strengthens the earlier source classification by using the actual rewrite
derivation to retain the complement of each redex. The Lean transition
system, original BPerm pair, stability/copyability premises, terms and
Figure~16 translation are unchanged. All contexts are allowed in this
BPerm result; the eventual term observations impose their original
binding-free restriction separately. The remaining ten template families,
correspondence work, semantic substitution/weakening and Figure~22 chain are
not consequences already established by this milestone.

Let $S$ denote the original BPerm specimen system. For arbitrary cost orders
$Q,P,R$, the final results are
\[
 S.\mathsf{Robust}(Q,P,R),\qquad
 S^{\mathrm{op}}.\mathsf{Robust}(Q,P,R).
\]
A matching case uses $k=0$ additional left steps and $l=1$ right step.
Its cost side condition is just $Q(1,1)$, hence reflexivity suffices.
The residual is first proved with equality on representation classes;
reflexive refinement then supplies the requested $P,R$ comparisons.
The exact-cost instance therefore agrees with the equality row of Table~3.

\lean{Copying/FullBPermRobustness}{32}{59}{template_robust}
\lean{Copying/FullBPermRobustness}{63}{91}{template_robust_converse}

\subsection{Atomic contributions and actual edgewise replacements}
\termdef{atomic-bperm-contribution}
An atomic contribution is a pair $A,B$ justified by one of three substantive
constructions: a recursive isomorphism; the original shallow BPerm relation;
or two boxes whose bodies have an original recursive filling witness.
Two administrative constructors compose endpoint isomorphisms and transport
equal interface lists. They change neither a rule nor an ordered port list.
Induction converts each atomic certificate into a genuine
\texttt{FillingWitness}. Conversely, inspecting an original edge certificate
extracts such an atom: retain its corolla isomorphism, its actual shallow
pair, or its recursively filled box body and endpoint corolla maps.

\lean{Copying/AtomicFillings}{12}{21}{AtomicPair}
\lean{Copying/AtomicFillings}{25}{34}{AtomicPair.witness}
\lean{Copying/AtomicFillings}{67}{85}{atomic_of_edge}

If the source has only primitive outer labels, an atomic contribution is
rigid. A shallow BPerm contribution and a boxed contribution both contain a
box edge, contradicting that hypothesis. The remaining isomorphism case is
immediate; the administrative cases transport the same argument. This proof
is about primitive footprints, not an assumption that every BPerm pair is
an ordered isomorphism.

\lean{Copying/AtomicFillings}{37}{63}{AtomicPair.primitive_rigid}

\termdef{bperm-decoration}
A decoration $D$ of $G$ assigns to every actual edge $e$ a replacement
$D_e$, with an atomic certificate from the labelled corolla of $e$ to $D_e$.
Its graph is the simultaneous expansion $G[D]$. The identity decoration
assigns the original corollas; collapsing these gives an explicit identity
isomorphism. Expanding the atomic filling witnesses and composing that
collapse proves that $G$ and $G[D]$ have an original filling witness whenever
$G$ is focus-free. For a primitive footprint every component is rigid;
the component isomorphisms assemble to $G\cong G[D]$. Its vertex map is
exactly the host inclusion, an essential fact for matching the active token.

\lean{Copying/DecoratedFillings}{9}{12}{Decoration}
\lean{Copying/DecoratedFillings}{17}{17}{Decoration.graph}
\lean{Copying/DecoratedFillings}{19}{21}{Decoration.identity}
\lean{Copying/DecoratedFillings}{23}{24}{Decoration.identityIso}
\lean{Copying/DecoratedFillings}{27}{29}{Decoration.witness}
\lean{Copying/DecoratedFillings}{32}{34}{Decoration.primitiveIso}
\lean{Copying/DecoratedFillings}{36}{38}{Decoration.primitiveIso_vertex}

\subsection{Normalising every original source specimen}
A recursive isomorphism pulls an ordered corolla certificate back along its
edge bijection. Consequently the left member of every shallow BPerm pair,
and then every original retained, shallow or boxed filling component,
is an ordered corolla. Its edge may have a different label from the common
context's hole, but its input and output signatures agree with that hole.
These facts are extracted from certificates, rather than imposed on specimens.

\lean{Copying/CorollaNormalisation}{11}{23}{LabelCopy.ofIso}
\lean{Copying/CorollaNormalisation}{27}{33}{UpToIso.left_copy}
\lean{Copying/CorollaNormalisation}{36}{43}{leftCopy}

\termdef{bperm-source-normalisation}
For a common context $C$ with filling families $L,R$, choose the actual
left label $a_e$ and its corolla certificate at each edge. Relabel $C$ by
$a_e$, keeping its vertices, edges and incidences; call the result $A$.
Collapsing $C[L]$ gives an explicit isomorphism $C[L]\cong A$.
Reinterface each $R_e$ by the proved signature equalities to obtain a
decoration of $A$. Its atomic evidence is the extracted original edge
certificate, with the same equality transports. Thus no shallow permutation
is replaced by an ordered isomorphism.

\lean{Copying/CorollaNormalisation}{49}{49}{SourceNormalisation.labels}
\lean{Copying/CorollaNormalisation}{50}{50}{SourceNormalisation.copies}
\lean{Copying/CorollaNormalisation}{53}{55}{SourceNormalisation.source}
\lean{Copying/CorollaNormalisation}{57}{59}{SourceNormalisation.leftIso}
\lean{Copying/CorollaNormalisation}{63}{64}{SourceNormalisation.right}
\lean{Copying/CorollaNormalisation}{66}{72}{SourceNormalisation.decoration}

The right comparison needs an expansion isomorphism whose skeleton labels
may differ but whose corresponding signatures are equal. A bijective map
is made into a bijection by choosing inverses; dependent sums of edge and
component bijections then give the expanded edge carrier. Expanded vertices
are the sum of the host vertices and component interiors. The host and
inserted-port equations verify their incidence maps, including the equality
casts on component boundaries. Mapping each source and target list proves
the resulting expansion isomorphism.

\lean{Graph/ExpansionTransport}{10}{15}{Bijection.ofBijective}
\lean{Graph/ExpansionTransport}{18}{30}{Bijection.sigmaMap}
\lean{Graph/ExpansionTransport}{43}{44}{Net.ExpansionTransport.vertices}
\lean{Graph/ExpansionTransport}{46}{47}{Net.ExpansionTransport.edges}
\lean{Graph/ExpansionTransport}{49}{51}{Net.ExpansionTransport.map_host}
\lean{Graph/ExpansionTransport}{54}{67}{Net.ExpansionTransport.map_insert}
\lean{Graph/ExpansionTransport}{71}{94}{Net.ExpansionTransport.iso}

Apply this construction to the identity incidence map from the relabelled
skeleton to $C$ and to each double interface-cast round trip. It gives
$A[D]\cong C[R]$. Both source-normalisation isomorphisms have explicit
host-token equations. They identify the actual coordinates used by a rewrite,
not just the unpointed graphs.

\lean{Copying/CorollaNormalisation}{75}{81}{SourceNormalisation.skeleton}
\lean{Copying/CorollaNormalisation}{83}{92}{SourceNormalisation.rightIso}
\lean{Copying/CorollaNormalisation}{94}{95}{SourceNormalisation.left_token}
\lean{Copying/CorollaNormalisation}{97}{99}{SourceNormalisation.right_token}

\subsection{Retaining complements and residual token positions}
\termdef{pointed-residual-filling}
A pointed filling from $(G,v)$ to $(H,w)$ consists of an original filling
witness, a vertex of its common context, and equations saying that the two
endpoint maps send that vertex to $v,w$. Its token is in search mode after
a rewrite. Boundary inputs give such certificates immediately. A reflexive
certificate works for any focus-free graph and any internal or boundary
vertex. Composing endpoint isomorphisms transports both equations, by
composition of vertex maps.

\lean{Copying/DecoratedFillings}{42}{47}{PointedFilling}
\lean{Copying/DecoratedFillings}{51}{52}{PointedFilling.input}
\lean{Copying/DecoratedFillings}{55}{56}{PointedFilling.refl}
\lean{Copying/DecoratedFillings}{59}{75}{PointedFilling.congr}

Simultaneous expansion of component witnesses can resume inside a selected
component's common context. The resulting point is its inserted vertex.
For an internal vertex the coordinate equality is definitional; for an input
or output it follows from the expansion host-position laws. A pointed closed
certificate yields the original equality-up-to quasi-specimen: pull the
term typing back through the endpoint isomorphism and host inclusion, then
apply the original quasi-specimen construction. Rootedness remains handled
by that definition.

\lean{Copying/PointedFillingAssembly}{10}{27}{PointedFilling.expandAt}
\lean{Copying/PointedFillingAssembly}{31}{44}{PointedFilling.quasi}

A dependent replacement family changes exactly one assigned graph. At the
selected edge it is the replacement; at every other edge it is the old
family member. These two equations let us combine a selected actual rewrite,
its pointed residual, and all independently supplied surrounding residuals.
The assembled right rewrite uses the original expansion constructor.
The assembled residual follows the selected common-context point, so the
resumed token need not be an outer boundary input.

\lean{Copying/PointedExpansion}{10}{18}{PointedFilling.replace}
\lean{Copying/PointedExpansion}{20}{23}{PointedFilling.replace_selected}
\lean{Copying/PointedExpansion}{25}{28}{PointedFilling.replace_other}
\lean{Copying/PointedExpansion}{32}{69}{PointedFilling.assemble}

\termdef{expanded-rewrite-presentation}
For the induction, an expanded rewrite has the original local, recursive
renaming and simultaneous-expansion constructors. Every such derivation is
an original rewrite. Ordinary plugging is derivable: fill all unselected
edges with identity corollas, peel the selected component out, and collapse
the remaining corollas. The selected inserted-port equation follows by
cancelling the peeling isomorphism with its inverse. Renaming the expanded
step along these two isomorphisms gives the original plug step at its exact
source and target tokens. Induction on the original rewrite therefore gives
the converse presentation theorem. This is an equivalence of presentations,
not a change to the machine's transition relation.

\lean{Copying/ExpandedRewrite}{9}{21}{Combined.ExpandedRewrite}
\lean{Copying/ExpandedRewrite}{25}{29}{Combined.ExpandedRewrite.actual}
\lean{Copying/ExpandedRewrite}{32}{41}{Combined.ExpandedRewrite.peel}
\lean{Copying/ExpandedRewrite}{43}{66}{Combined.ExpandedRewrite.peel_insert}
\lean{Copying/ExpandedRewrite}{69}{93}{Combined.ExpandedRewrite.plug}
\lean{Copying/ExpandedRewrite}{97}{103}{Combined.Rewrite.expanded}

A decoration pulls back through a recursive source isomorphism by its edge
bijection and the corresponding signature equalities. Expansion transport
provides the whole-graph comparison, with its host-token law. For an expanded
source, restrict a decoration to each already-present component. Expansion
associativity regroups the resulting graph and proves the equation for a
token inserted into the selected component. Thus the induction carries the
actual complement supplied by the rewrite derivation.

\lean{Copying/DecorationTransport}{12}{20}{Decoration.pull}
\lean{Copying/DecorationTransport}{22}{29}{Decoration.pullIso}
\lean{Copying/DecorationTransport}{31}{33}{Decoration.pullIso_host}
\lean{Copying/DecorationTransport}{41}{43}{Decoration.component}
\lean{Copying/DecorationTransport}{45}{46}{Decoration.groupedIso}
\lean{Copying/DecorationTransport}{49}{54}{Decoration.groupedIso_insert}

The local transfer proposition explicitly asks for an actual right rewrite
and a pointed original residual. Assuming that proposition temporarily,
induction on expanded rewrites proves whole-graph transfer: use the local
case, pull back in the renaming case, and restrict/regroup/assemble in the
expansion case. The equivalent rewrite presentation removes the temporary
presentation choice. Primitive-only sources discharge the local obligation
by the rigidity isomorphism and the unchanged reduct's reflexive certificate.
The later theorem \texttt{local\_transport} discharges the temporary local
hypothesis exhaustively; it is not a premise of full robustness.

\lean{Copying/WholeRewriteTransport}{11}{14}{LocalTransport}
\lean{Copying/WholeRewriteTransport}{17}{56}{transport_expanded}
\lean{Copying/WholeRewriteTransport}{59}{63}{transport_rewrite}
\lean{Copying/WholeRewriteTransport}{67}{76}{primitive_transport}

\subsection{Canonical box cases and unchanged local operations}
\termdef{canonical-bperm-box-view}
An atomic normal form collects all endpoint representation changes. It is
an isomorphism, a literal shallow pair with its two endpoint maps, or a boxed
body filling with its endpoint maps. Interface equality transports are
retained explicitly. Congruence composes the maps; interface transport
eliminates equalities; inspection of the original shallow witness gives the
literal pair; induction normalises an arbitrary atomic certificate.

\lean{Copying/AtomicNormalForm}{10}{24}{AtomicNormal}
\lean{Copying/AtomicNormalForm}{28}{33}{AtomicNormal.congr}
\lean{Copying/AtomicNormalForm}{35}{38}{AtomicNormal.interfaces}
\lean{Copying/AtomicNormalForm}{41}{48}{UpToIso.normal}
\lean{Copying/AtomicNormalForm}{51}{60}{AtomicPair.normal}

When the source is a particular box $\Box B$, its view is retained, deep
with a witness $B\leadsto C$, or shallow with a boundary isomorphism
$B\to C$ and its exact free-output action. The proof compares the source
box label with the normal form's box label. Recursive label equivalence
forces the same binder and free-interface metadata and supplies the body
isomorphism. In the deep case compose that isomorphism with the body witness;
in the shallow case compose it with the inverse boundary permutation.
The right endpoint map still includes the outer inverse compensation.

\lean{Copying/AtomicBoxView}{8}{15}{BoxView}
\lean{Copying/AtomicBoxView}{19}{45}{AtomicPair.boxView}

Translated lambda boxes use $k+1$ in the body interface, whereas the generic
box constructor uses $1+k$. The cast round-trip lemma proves that the two
equality transports cancel. Undoing the source-body cast transports a
boundary isomorphism and its output action; the proof is equality elimination.
The lambda view then has the same three cases at the original translation's
interface. Uniqueness of the pair's inner action identifies it with the
one-bound-wire action required by the existing opening diagram.

\lean{Copying/LambdaBoxView}{9}{11}{cast_roundtrip}
\lean{Copying/LambdaBoxView}{13}{15}{undoBodyCast}
\lean{Copying/LambdaBoxView}{17}{21}{undoBodyCast_action}
\lean{Copying/LambdaBoxView}{23}{29}{LambdaView}
\lean{Copying/LambdaBoxView}{32}{58}{AtomicPair.lambdaView}

Decorations distribute over the actual graph constructors. A decorated
corolla collapses to its selected filling by peeling and the corolla plug
unit. Serial left/right decorations are restrictions to the original two
slots, and their comparison is expansion associativity. Tensor left/right
decorations restrict to the two branch tags; parallel expansion gives their
comparison. Equality-interface transport is undone by equality elimination.
The tensor congruence used below assembles the two component isomorphisms.

\lean{Copying/DecorationAssembly}{9}{12}{Decoration.corollaIso}
\lean{Copying/DecorationAssembly}{14}{16}{Decoration.serialLeft}
\lean{Copying/DecorationAssembly}{18}{20}{Decoration.serialRight}
\lean{Copying/DecorationAssembly}{22}{28}{Decoration.serialIso}
\lean{Copying/DecorationAssembly}{30}{33}{Decoration.tensorLeft}
\lean{Copying/DecorationAssembly}{35}{38}{Decoration.tensorRight}
\lean{Copying/DecorationAssembly}{40}{47}{Decoration.tensorIso}
\lean{Copying/DecorationAssembly}{50}{53}{Decoration.unboundary}
\lean{Copying/BetaDecoration}{9}{11}{tensor_congr}

For beta, \texttt{withBox} is the literal application graph with just its
lambda box varied. Its congruence follows from serial and tensor congruence.
The decoration's box slot is obtained by selecting the application operands,
the lambda operand and its box. All other local contributions are primitive:
the application edge, lambda edge and the mandatory stable argument.
Their rigidity isomorphisms therefore identify the decorated source with
this actual variable-box graph. Moving the shallow inverse compensation
through lambda, tensor and application gives the previously proved
compensated source, with the same ordered free and argument output blocks.

\lean{Copying/BetaDecoration}{17}{19}{BetaDecoration.withBox}
\lean{Copying/BetaDecoration}{21}{23}{BetaDecoration.withBox_congr}
\lean{Copying/BetaDecoration}{27}{27}{BetaDecoration.boxes}
\lean{Copying/BetaDecoration}{30}{37}{BetaDecoration.sourceIso}
\lean{Copying/BetaDecoration}{40}{49}{BetaDecoration.compensated}

Now inspect the lambda view. The retained case uses the original beta step
and a reflexive residual. The deep case opens the paired bodies, preserving
the recursively filled witness through the original depth increase and
bound-output splice. The shallow case uses the original interface-permuted
beta step and its compensated successor isomorphism, hence a genuine
residual witness. Focus freedom of the body is extracted from the source;
stability of the argument is the original rule premise, not a conclusion
inferred from linear semantics. Every case resumes at input zero of the
local reduct, ready for arbitrary surrounding assembly.

\lean{Copying/BetaDecoration}{53}{76}{BetaDecoration.transport}

For copying a lambda payload, select its box slot and collapse only the
primitive contraction tree and lambda head. The three box views give a
retained copy step, a deep copy step whose two branches both contain the
body filling, or the shallow compensated copy step with separately certified
residual branches. The contraction tree and the output distributors remain
the original constructions. For the complete copyable datatype, instance,
constant, application, unary and binary payloads are primitive. The lambda
case is the preceding proof. Recursive renaming pulls the decoration back
along the whole copying source isomorphism and transports the entire
successor and pointed residual along the copying successor isomorphism.
Thus every constructor of copyability is handled.

\lean{Copying/CopyDecoration}{13}{13}{CopyDecoration.boxes}
\lean{Copying/CopyDecoration}{15}{21}{CopyDecoration.sourceIso}
\lean{Copying/CopyDecoration}{23}{49}{CopyDecoration.lambda_transport}
\lean{Copying/CopyDecoration}{53}{85}{CopyDecoration.transport}

Exhaustive case analysis on the original local-rule datatype now proves local
transfer: allocation, dereference, assignment, both equality tests and both
arithmetic families use their primitive footprints; beta and copying use
the proved diagrams above. Substitution into the whole-graph induction removes
its local hypothesis. Finally the source-normalisation isomorphisms transport
an arbitrary original $C[L]$ rewrite to a decoration and the matched rewrite
back to $C[R]$, preserving the actual host token on both sides.

\lean{Copying/BPermForwardTransport}{10}{21}{local_transport}
\lean{Copying/BPermForwardTransport}{24}{27}{rewrite_transport}
\lean{Copying/BPermForwardTransport}{31}{52}{filling_rewrite}

\subsection{Converse transfer without assuming symmetry of the raw pair}
\termdef{reversed-bperm-filling}
First, an external typed output permutation acts on a filling witness by
putting it inside an actual permuted corolla context. Expanding its component
witness and using the single-expansion and permuted-corolla plug isomorphisms
identifies both endpoints with the externally permuted graphs.

To reverse a shallow pair, start with the inverse pair on its unpermuted
body, apply the original outer inverse permutation to both endpoints, and
cancel the resulting double outer permutation. Compose the original endpoint
isomorphisms. This constructs a genuine reversed filling; it does not assert
that the directed raw BPerm relation itself is symmetric.

\lean{Copying/PermutedFilling}{10}{38}{FillingWitness.permuted}
\lean{Copying/PermutedFilling}{46}{60}{UpToIso.reverseWitness}

A reversal stores a reversed filling for every original common-context edge.
Expanding these components and exchanging the original endpoint isomorphisms
gives a reversed whole witness. The new common point is the host image of
the old point, so its endpoint coordinates are exactly exchanged. Construct
these components by recursion on box depth: retained corollas reverse by
isomorphism, shallow pairs use the preceding construction, and boxed pairs
recursively reverse their actual body witness and rebuild the two box
corollas. The recursion decreases the depth at each boxed call. Applying
this construction to a pointed filling proves its symmetry while retaining
both exact token coordinates.

\lean{Copying/ReverseFillings}{9}{10}{Reversal}
\lean{Copying/ReverseFillings}{15}{16}{Reversal.witness}
\lean{Copying/ReverseFillings}{19}{24}{Reversal.pointed}
\lean{Copying/ReverseFillings}{29}{45}{reverseFilling}
\lean{Copying/ReverseFillings}{48}{52}{PointedFilling.symm}

Forward pointed rewrite transport pulls the supplied step back through the
left endpoint isomorphism, applies the completed common-context rewrite
theorem, and pushes the matched step through the right endpoint isomorphism.
The point equations cancel the inverse maps. For converse transport, reverse
the source certificate, apply this forward theorem to the actual right step,
and reverse the resulting residual. The residual is thus returned in the
original left-to-right template orientation.

\lean{Copying/PointedRewriteTransport}{10}{32}{PointedFilling.rewrite}
\lean{Copying/PointedRewriteTransport}{36}{40}{PointedFilling.rewrite_converse}

\subsection{Robustness, exact execution and regression checks}
Every original specimen yields a pointed filling directly from its existing
context, edge certificates and host token. A representation-class step from
rewrite mode reflects to an actual finite-depth rewrite: class reflection
preserves the allowance, and the search alternative is impossible in rewrite
mode. Apply forward or converse pointed transport to that actual rewrite.
Rewrite preservation proves both residual graphs hole-free; endpoint typing
supplies the resumed term token. The pointed residual therefore supplies the
original equality-up-to quasi-specimen. Projection gives matching class
steps; machine determinism identifies any independently presented successor
with the constructed one. This proves the second robustness disjunct with
$k=0,l=1$, in both directions, for every specimen and cost-order triple.

\lean{Copying/FullBPermRobustness}{10}{16}{specimenPointed}
\lean{Copying/FullBPermRobustness}{20}{28}{class_rewrite}

The initial search coordinates of any closed paired filling are rooted and
remote. Its actual filling witness therefore supplies the original lifting.
Apply the already proved sufficiency-of-robustness theorem in both directions,
using BPerm input safety and the exact reasonable cost triple. This gives
exact quantitative equivalence of the actual unbounded initial states.
Unfolding the two refinements proves equality of the successful-run spectra:
for each $n$, one graph terminates in $n$ transitions if and only if the other
does. No termination assumption on callbacks or box bodies is needed.
This is a theorem about arbitrary certified closed graph fillings; it does
not by itself complete the remaining structural derivations or Figure~22.

\lean{Copying/BPermExactExecution}{10}{23}{filling_initial_lifting}
\lean{Copying/BPermExactExecution}{27}{35}{filling_exact_equiv}
\lean{Copying/BPermExactExecution}{39}{52}{filling_termination_count}

The new regressions instantiate universal transfer in a host whose active
position is internal and whose other component contains separately changed
shallow and nested BPerm contributions. One performs arithmetic; the other
opens a deeply related beta body. Both conclusions contain an actual right
whole-graph rewrite and a constructed original residual, improving the earlier
opposite-match-only regression. The final check records that the chosen host
token is not an input boundary coordinate.

\lean{Copying/FullBPermExamples}{11}{25}{FullExamples.arithmetic_mixed_host}
\lean{Copying/FullBPermExamples}{29}{44}{FullExamples.beta_mixed_host}
\lean{Copying/FullBPermExamples}{47}{49}{FullExamples.host_token_not_input}
